% ==============================================================================
% 前言
% ==============================================================================
\clearpage
\phantomsection
\thispagestyle{plain}
\markboth{前言}{前言}
\addcontentsline{toc}{section}{前言}

\begin{center}
  {\Large\sffamily\bfseries\color{accentblue} 前\qquad 言}
\end{center}
\vspace{4mm}

各种各样的曲线，形状千姿百态，性质丰富多彩，应用举不胜举。无论从哪一方面看，用“形形色色”这四个字来形容曲线，都显得十分恰当。这本以《形形色色的曲线》为题的小册子，就是试图在中学数学知识的基础上，谈谈各种著名的、美丽的和有用的特殊曲线。读者看这本书，一方面可以把刚刚学过的解析几何与微积分知识用起来，另一方面可以增长知识，弄清楚一些有趣事实的数学原理。

你知道电影里是用什么方法来产生光芒四射的效果吗？那里就是利用了某一种曲线的性质，而且这种曲线还有许多奇妙的性质，以至于一位大名鼎鼎的数学家特地在遗嘱里吩咐要把它刻在自己的墓碑上。你一定知道两个物体从不同的高度一起下落，下面的物体先落地，然而你是不是知道，两个物体同时从不同的高度沿着某一种曲线下滑，却能同时滑到终端？……诸如这类，本书当然都要介绍。不过，为了让读者能从阅读中获得最大的乐趣，在前言里还是少说为宜。一面阅读，一面思考，一面联想，或许更有意思些。

为了把一些专门知识通俗地叙述出来，本书难免出现各种缺点和错误，恳请读者多多批评指正。

\vspace{6mm}
\begin{flushright}
作\quad 者\hspace*{18mm}\\[2mm]
1983 年 8 月 1 日于扬州\hspace*{4mm}
\end{flushright}
\clearpage
